\section{描述性统计与来源结构演化}

表 \ref{tab:descriptive} 给出多产品面板的基本描述统计。四类产品均覆盖 17 个年份，但出口国数量和贸易规模存在明显差异。854231、854232 和 854239 的出口国数量高于 848620，说明集成电路类产品的全球贸易来源更广；但更广的来源并不必然意味着风险更低，因为部分产品仍可能集中于少数核心生产网络节点。

\begin{table}[H]
\centering
\scriptsize
\caption{多产品面板描述统计}
\label{tab:descriptive}
\begin{tabular}{lrrrrr}
\toprule
指标 & 全部产品 & 848620 & 854231 & 854232 & 854239 \\
\midrule
观测值 & 10608 & 1020 & 3366 & 2737 & 3485 \\
正向贸易记录 & 5619 & 585 & 1789 & 1251 & 1994 \\
出口国-产品组合 & 624 & 60 & 198 & 161 & 205 \\
年份数 & 17 & 17 & 17 & 17 & 17 \\
平均进口额（千美元） & 246456.38 & 172161.32 & 321346.33 & 261591.24 & 183982.15 \\
平均进口份额 & 0.006410 & 0.016667 & 0.005051 & 0.006211 & 0.004878 \\
\bottomrule
\end{tabular}
\end{table}


图 \ref{fig:total-imports} 展示四类产品的总进口规模变化。不同产品的规模波动差异较大，说明中国半导体相关进口不仅受政策因素影响，也与产业周期和国内需求变化有关。

\begin{figure}[H]
    \centering
    \includegraphics[width=0.92\textwidth]{total_imports_by_product_v02.png}
    \caption{中国半导体相关产品总进口额变化}
    \label{fig:total-imports}
\end{figure}

从美国进口份额看，848620 的美国份额从 2017 年 27.92\% 降至 2024 年 9.88\%，854231 从 12.47\% 降至 9.72\%，854232 从 0.95\% 降至 0.14\%。854239 则从 2.82\% 升至 3.94\%，表明集成电路类产品并非完全同向变化。

\begin{figure}[H]
    \centering
    \includegraphics[width=0.92\textwidth]{usa_import_share_by_product_v02.png}
    \caption{美国来源进口份额变化}
    \label{fig:usa-share}
\end{figure}

2024 年主要来源国结构显示，设备类产品更多依赖日本、荷兰、新加坡和美国等设备强国；集成电路类产品则更多体现韩国、马来西亚、越南和 Other Asia, nes 等亚洲生产网络节点的作用。

\begin{figure}[H]
    \centering
    \includegraphics[width=0.92\textwidth]{top_2024_source_shares_by_product_v02.png}
    \caption{2024年主要来源国份额}
    \label{fig:top-sources}
\end{figure}

HHI 指标进一步说明，来源替代不等同于来源分散。部分产品在美国份额下降的同时，仍集中于少数主要来源国，因而需要综合指标评估供应链风险。

\begin{figure}[H]
    \centering
    \includegraphics[width=0.92\textwidth]{source_hhi_by_product_v02.png}
    \caption{进口来源集中度 HHI}
    \label{fig:hhi}
\end{figure}
